% %%%%%%%%%%%%%%%%%%%%%%%%%%%%%%%%%%%%%%%%%%%%%%%%%%%%%%%%%%%%%%%%%%%%%
% %% SpDCache %%
% Introduction
% %%%%%%%%%%%%%%%%%%%%%%%%%%%%%%%%%%%%%%%%%%%%%%%%%%%%%%%%%%%%%%%%%%%%%

\section{Introduction}
\label{sec:spdc:intro}

\lettrine{I}{n} the previous chapters, we analyzed the causes and effects of irregular memory accesses in sparse computations and identified that the \acrfull{op} algorithm of \acrfull{spmv} ($A \times b = c$) has potential for high performance and optimization.
The \acrshort{op} algorithm generates irregular updates on the output vector $c$, called ``random'' \glspl{reduction}, as each non-zero element of the input matrix $A$ produces updates to distinct, and often distant, location in the output vector $c$.
These ``random'' \glspl{reduction} render stress the memory system.

Conventional cache hierarchies are designed for regular, predictable access streams that exploit spatial and temporal locality.
In contrast, \acrshort{op} \acrshort{spmv} exhibits poor locality, leading to low cache utilization, frequent \glspl{eviction}, and thus requiring more memory bandwidth for the output vector $c$.

Previous works, such as Coup~\cite{Zhang2015_Coup}, PHI~\cite{Mukkara2019_PHI}, SortCache~\cite{Srikanth2021_SortCache}, and ASA~\cite{Zhang2022_ASA}, have proposed \gls{reduction}-aware cache mechanisms to alleviate irregularity by improving accumulation and storage efficiency.
However, these solutions do not exploit the structural properties of the processed matrices, and therefore miss optimization opportunities when the matrix has some regularity.
This observation motivates our focus on \gls{banded} sparse matrices, which combine overall sparsity with a dense region near the diagonal.
\Gls{banded} matrices are very common and serve as an initial example to study how dense and sparse regions can be leveraged to enhance performance.

This chapter begins by reviewing the main concepts of generic data-caches, introducing certain terminology and mechanisms which are required to discuss cache behavior under \gls{irrAccess} conditions.
It then presents a \gls{reduction}-aware cache architecture, designed to efficiently support ``random'' \glspl{reduction} for the \acrshort{op} \acrshort{spmv} kernel.
The proposed design, \acrlong{spdc} (presented in Paper~\ref{publ:SpDC}), extends traditional caching models with a region-based organization and native \gls{reduction} support.
Finally, we evaluate its performance through high-level modeling on \gls{banded} and non-\gls{banded} matrices, highlighting the influence of data regularity on cache efficiency.

\clearpage